\documentclass[a4paper,twoside]{scrartcl}

\usepackage{graphicx,color}
\graphicspath{{images/}{../htdocs/images/header/}}

\usepackage{mparhack}

% \usepackage[a4paper]{geometry}
% \geometry{%
% twoside,%
% %height=.7\paperheight,%
% %width=.7\paperwidth,%
% top=.1\paperheight,%
% bottom=.2\paperheight,%
% left=.1\paperwidth,%
% right=.2\paperwidth%
% %bindingoffset
% }

\usepackage{fancyhdr}
\pagestyle{fancy}
\fancyhead{}
\fancyfoot{}
\fancyfoot[LE,RO]{\iffloatpage{}{\thepage}}
\renewcommand{\headrulewidth}{0pt}
\renewcommand{\footrulewidth}{0pt}
\setlength{\oddsidemargin}{.1\paperwidth}
\setlength{\evensidemargin}{.2\paperwidth}
\setlength{\headheight}{.1\paperheight}
\setlength{\headsep}{0pt}
\setlength{\topmargin}{0pt}
\setlength{\voffset}{-1in}
\setlength{\hoffset}{-1in}
\setlength{\textwidth}{.7\paperwidth}
\setlength{\textheight}{.7\paperheight}
\setlength{\marginparwidth}{.1\paperwidth}

\newcommand{\nicesection}[2]{%
\cleardoublepage
\fbox{\includegraphics[width=\linewidth]{#1}}%
\vspace*{-1em}%
\section{\hfill #2}
\hrule
\vspace*{\baselineskip}%
}

\fboxsep0pt
\setlength{\parindent}{0pt}

\newcommand{\todo}[1]{{\color{red}\bf TODO: #1}}
\newcommand{\comment}[1]{}
\definecolor{codecol}{rgb}{0.1,0.2,0.5}
\newcommand{\code}[1]{\texttt{\color{codecol}#1}}

\frenchspacing
% \usepackage[T1]{fontenc}
% \usepackage[condensed,math]{iwona}

%\usepackage[T1]{fontenc}
%\usepackage[urw-garamond]{mathdesign}

\usepackage[T1]{fontenc}
\usepackage[scaled]{helvet}
\renewcommand{\familydefault}{\sfdefault}

%\setkomafont{sectioning}{\sfdefault}


\title{Darktable Programmer's Guide}
\author{hanatos}

\begin{document}

\fbox{\includegraphics[width=\linewidth]{header12}}%
\vspace*{-1em}%
\section*{\hfill Darktable Programmer's Guide}
\hrule

\vspace*{4\baselineskip}

{\hfill version \input version.tex }

\thispagestyle{empty}

\newpage
\tableofcontents


\nicesection{header1}{Introduction}
\label{sec:introduction}


\nicesection{header2}{System Layout}

\resizebox{\linewidth}{!}{\input{graphs/system.pdftex_t}}

\subsection{Module Descriptions}

\begin{description}
  \item[gui] using gtk2 and cairo.
\end{description}
\begin{description}
  \item[control]
  \item[control, scheduler] using pthreads and a custom job queuing system.
\end{description}
\begin{description}
  \item[image]
  \item[image cache]
  \item[mipmap cache]
  \item[db]
  \item[imageio]
\end{description}
\begin{description}
  \item[library] this is the lighttable module.
  \item[film]
\end{description}
\begin{description}
  \item[develop] this implements the darktable.
  \item[pixelpipe]
  \item[image operation (iop) modules, aka plug-ins]
\end{description}

\newpage
\subsection{Image Loading Data Flow}

\paragraph{Related Modules}
\begin{itemize}
  \item imageio \code{dt\_imageio\_*}
  \item db (messily hidden in sqlite3 statements all over the place)
  \item image cache \code{dt\_image\_cache\_*}
  \item mip cache \code{dt\_mipmap\_cache\_*}
  \item image \code{dt\_image\_t}
\end{itemize}

\paragraph{in image cache}
\begin{itemize}
  \item \code{dt\_image\_import}

  db $\rightarrow$ image struct, trigger imageio disk $\rightarrow$ image entry in db if miss,
                metadata through libraw/magick
  
 
\end{itemize}

\paragraph{in mip cache}
\begin{itemize}
  \item \code{dt\_image\_get}

      db $\rightarrow$ mip[0--4], trigger imageio disk $\rightarrow$ mip4 and mip4 $\rightarrow$ mip[0-3] if miss.

  
      db $\rightarrow$ mipf, trigger mip4 $\rightarrow$ mipf or pixels $\rightarrow$ mipf if miss.
\end{itemize}

\paragraph{in imageio}
\begin{itemize}
  \item \code{dt\_imageio\_open\_preview}

    disk $\rightarrow$ mip[0--4] entries in db (load thumbnail).
    \begin{itemize}
      \item[\todo{}]  \code{dt\_imageio\_open\_ldr\_preview}
      \item \code{dt\_imageio\_open\_raw\_preview}
    \end{itemize}
  \item mip4 $\rightarrow$ mipf (by guessed reverse gamma) \code{dt\_image\_preview\_to\_raw}.
  \item full pixels $\rightarrow$ mipf (downscaling) \code{dt\_image\_raw\_to\_preview}.
  \item mip4 $\rightarrow$ mip[0--3] \code{dt\_image\_update\_mipmaps}.
\end{itemize}

image import: only load \code{dt\_image\_t} and mip[0--4] to database.

\code{dt\_image\_get}: check mipmap cache, check database, else \code{dt\_imageio\_open[\_preview]}.

\newpage
\subsection{Cache Interfaces}

\subsubsection{Image Cache}

\code{dt\_image\_cache\_t}

\subsubsection{Mipmap Cache}

\code{dt\_image\_t}

\begin{description}
  \item{\code{dt\_image\_get}} get buffer or, if missed, a lower resolution (launch job in bg for correct resolution)
  \item{\code{dt\_image\_load}} Load buffer in this thread, return exactly this resolution.
    No locking is performed, so be sure to acquire the lock using \code{dt\_image\_alloc} before.
\end{description}

\subsubsection{Homebrew Pixelpipe Cache}

\newpage
\subsection{Plug-in Interface}

For examples, have a look at \code{src/iop/*.\{h,c\}}, and also see the method declarations in \code{src/common/imageop.h}.

Pixel processing is done in pipelines (\code{src/develop/pixelpipe\_hb.c}) with fixed order of the modules,
observe the current status with e.g.

\code{grep priority src/iop/*.c | sort -t '=' -k 2}.

At the time of this writing, the main pixel pipeline stack is organised as follows:

\medskip

\begin{tabular}{lrl}
module name & priority & input color format \\
\hline
rawimport      &  100  &  - \\ 
exposure       &  150  &  fff input color space \\
temperature    &  200  &  fff input color space \\
highlights     &  250  &  fff input color space \\
basecurve      &  260  &  fff input color space \\
profile\_gamma &  299  &  fff input color space \\
colorin        &  300  &  fff input color space \\
colortransfer  &  350  &  fff $L^*a^*b^*$ \\
equalizer      &  500  &  fff $L^*a^*b^*$ \\
monochrome     &  550  &  fff $L^*a^*b^*$ \\
tonecurve      &  700  &  fff $L^*a^*b^*$ \\
colorzones     &  750  &  fff $L^*a^*b^*$ \\
colorcorrection&  800  &  fff $L^*a^*b^*$ \\
colorout       &  900  &  fff $L^*a^*b^*$ \\
lens           &  940  &  fff output color space \\
clipping       &  950  &  fff output color space \\
grain          &  965  &  fff output color space \\
clahe          &  966  &  fff output color space \\
velvia         &  967  &  fff output color space \\
splittoning    &  968  &  fff output color space \\
vignette       &  970  &  fff output color space \\
sharpen        &  990  &  fff output color space \\
gamma          & 1000  &  fff output color space \\
               &   -   &  ccc output color space \\
\end{tabular}

\medskip

where fff is $3\times32$ bit float, ccc is $3\times8$ bit int, and the input
color space depends on your camera or type of input data, and the output color space will
be sRGB (including gamma correction, i.e.\ it is non-linear) most of the time.

$L^*a^*b^*$ is a special variant of $L^*ab$, where $L^* \in [0, 100]$ and $a, b \in [-128, 128]$, but $a^* = a/L^*$ and $b^* = b/L^*$.
This makes sure that changes in $L^*$ transparently result in the correct adjustment in saturation:
$a_{out} = a_{in} \cdot L^*_{out}/L^*_{in}$ and
$b_{out} = b_{in} \cdot L^*_{out}/L^*_{in}$.


Each module has to implement a \code{process} function, which is responsible for processing of three types of
pipelines: preview, full, and export. Most of the time, this can be handled transparently, i.e.\ needs no special
handling in the module. The pipelines are as follows:

\begin{description}
  \item[preview] all of the image, but downsampled to low resolution for the navigation view, approximate full image statistics, and fast scrolling.
  \item[full] full resolution, but only the part that can currently be seen (the region of interest, roi, will be adjusted accordingly).
  \item[export] full image at full resolution, no gui data is inited for the module.
\end{description}

and can be identified by the pipeline struct, which carries a \code{pipe->type} flag. This struct can in turn be found
via the pipeline data piece which is passed to \code{process}: \code{piece->pipe}.

\paragraph{Regions of Interest}

A region of interest holds $x,y,w,h$, and scale information. The semantics are as follows: multiply the whole resolution image by
the scale value, then crop it to the window defined by $x, y, w, h$. The passed input buffer will be a contiguous block of values as defined by
the input color space, pixel by pixel and row by row. The output buffer should be in the same format, but the extent should match the
output region of interest.


\subsubsection{Most important Gui Callbacks}

\begin{description}
  \item \code{gui\_init} gui thread, init the widget (not called for export)
  \item \code{gui\_update} gui thread, make the gui match the modules parameters (not called for export)
  \item \code{commit\_params} core thread + gui thread in synch, copy module->params to module->data
  \item \code{process} core thread, execute the op on the input buffer
  \item \code{modify\_roi\_in}
  \item \code{modify\_roi\_out} core thread, optional, only needed if input and output buffers are of different size.
    For examples, see \code{src/iop/lens.c} and \code{src/iop/clipping.c}.
\end{description}

\subsubsection{Data Structures}

\begin{description}
  \item[\code{dt\_iop\_*\_params\_t module->\{params, default\_params, factory\_params\}}]

  This struct should only be modified by the gui thread. The three fields exist only once per module: \code{params}
  store the currently active parameters, \code{default\_params} are the default parameters (possibly adjusted by the
  user), and \code{factory\_params} are dt's factory defaults (automatically copied by the core after \code{init} has been
  called).

  \item[\code{dt\_iop\_*\_data\_t piece->data}]
  This struct lives in pixel pipeline domain, and thus exists once for the full and once for the preview pipeline.
  It should only be modified by the core thread (which calls \code{commit\_params} and \code{process}).

  \item[\code{dt\_iop\_*\_gui\_data\_t module->gui\_data}]
  You can use this struct to store your gui data (from the gui thread, obviously). It exists once per module and should be inited in \code{gui\_init} and
  cleaned up in \code{gui\_cleanup}. Note that gtk+ widgets are cleaned up implicitly while the widget is destroyed.

  \item[\code{dt\_iop\_*\_global\_data\_t module->data}]
  Core thread global data for all pipes (such as lensfun description database).
\end{description}


\subsubsection{A most useless Example}


\todo{script to generate as up-to-date minimal dt header set to compile external plugins}

\todo{a silly template module line by line: render a checker board on top of your input image.}

See \code{src/iop/useless.c}.

Change \code{src/iop/Makefile.am}:
\begin{verbatim}
libuseless_la_SOURCES=useless.c
dtplugins_LTLIBRARIES=libtonecurve.la [..] libcolortransfer.la libuseless.la
\end{verbatim}

\end{document}
